\section{Discussion and Summary}
\label{sec:discussion}

Perhaps the most useful aspect of considering system extension is that it elegantly unifies both parameter perturbations and application of harmonic forcings into a single framework. When viewed through the lens of an extended system a parameter perturbation is nothing but a harmonically forced system, albeit with a zero angular frequency and nonlinear effects. For the extension of the center subspace the only thing that is important is whether the extended eigenvalues are in resonance with the intrinsic critical eigenvalues of the system. In the previous sections parameter perturbations and harmonic forcings have been treated as separate cases however, this is primarily to illustrate the identical algorithmic treatment required for the extension of the tangent space in the two cases. 

Secondly, the system extension also illuminates the structure of the problem under forcing or parameter perturbation, with the extended directions effectively being the resolvent responses of the system. Of prime importance was the evaluation of the adjoint extended modes of the system that are essential in constructing the projection operators for the individual subspace directions. The extended adjoint modes have a special structure regardless of resonance, which was fruitfully utilized to show that, when constructing asymptotic center-manifold approximations of the extended system, all evaluations of the asymptotic vectors, $\AsECenM_{i_{1},\ldots,i_{k}}$, in fact only require the original linear system $\Lu$. This is an important aspect when considering approximations for large systems where specialized algorithms are often tailored to utilize the internal structure of $\Lu$. Despite the system extension, only the original linear system $\Lu$ is required for the calculations, which means that the specialized algorithms do not require any substantial modifications. In that sense the current approach provides the best of both worlds. It allows the calculation of center-manifold approximations of perturbed systems using the original unperturbed linear operators. At the same time it provides the theoretical perspective of understanding the perturbations as intrinsic modes of an extended system. 

It is also shown that the $\Khat$ matrix, which is the linear part of the reduced representation $\Param(\vecparv)$, is responsible for the cross coupling of vectors $\AsCenM_{i_{i},\ldots,i_{k}}$ at each order $k$ in the asymptotic center-manifold approximations. However, this matrix is upper triangular, at least in the normal form and graph-method approximation described in sections~\ref{sec:normal_form} and \ref{sec:graph_method}. The consequence of the upper triangular nature of $\Khat$ is that the linear system of equations obtained at each order for the solution of the asymptotic vectors $\AsCenM_{i_{i},\ldots,i_{k}}$ is lower triangular. Therefore at every order $k$, each vector can be solved for, one at a time, sequentially using back substitution. For large systems this is an important property since the coupled system solution can quickly become computationally infeasible. 

 Finally, the derived center-manifold approximation is applied to a model Ginzburg-Landau problem at the bifurcation point. The system is considered with perturbations of the diffusion coefficient $(\delta_{4})$ as well as with resonant and non-resonant external forcings. In all cases a very good agreeement is found between the non-linear system response and the response predicted by the calculated center-manifold approximations. 


%Furthermore, since zero angular frequencies are readily included into the extended system, this allows one to incorporate arbitrary steady external forcing terms into the center-manifold framework. 











